\begin{abstract}
Estimating metric camera displacement from a single low-cost inertial
measurement unit (IMU) is ill-conditioned, because acceleration does not
reveal the initial velocity and attitude errors leak gravity into
translation. We present a measurement framework for a handheld camera with an
endoscopic field of view in which inference uses only the 200-Hz six-axis IMU
stream, while chessboard poses serve solely for calibration, training
supervision and evaluation. The framework comprises a statistically
characterized planar reference, a digital twin fitted to the real recordings,
and PhysNet, a 0.42-M-parameter network that maps anchor-frame pre-integrated
IMU features onto a densely supervised velocity stream with a learned
stillness gate, a lever-arm estimate and per-prediction uncertainty.
Decomposing each 3-s displacement into an IMU-observable double integral and
an undetermined initial-velocity term explains the main empirical finding:
more than forty estimator variants share one error plateau, with predictions
shrunk toward zero. PhysNet lowers the synthetic-test error of an adapted
IMUNet by 24\% with one seventh of its parameters, and is statistically
indistinguishable from it on real data; all learned variants clearly beat
zero motion (42.2--43.6 versus 49.7~mm on a held-out recording). A protocol
study in the twin identifies the stop rate as the dominant acquisition
factor: one firm stop every 4~s lowers the error by 32\% relative to
pause-free motion. With gyroscope-integrated attitude and periodic position
anchoring, the translation stream attains a relative trajectory error of
80.6~mm against 119.2~mm for zero motion.
\end{abstract}

\begin{IEEEkeywords}
Camera motion, digital twin, inertial measurement unit, measurement
uncertainty, observability, pose measurement, simulation-to-real transfer.
\end{IEEEkeywords}
